% ---------------------------------------------------------------------------
% Table I -- the grid. Generated by experiments/make_grid_table.py; do not
% edit the numbers here, because there are none to edit. Regenerate instead.
%
% Spans both columns, so table* and not table. Needs \usepackage{booktabs}.
%
% The caption is ONE sentence and the note is a KEY, not an argument. ieeeconf
% sets table captions in small capitals, which is unreadable past a line, and
% the first draft ran to six. The note under the rule used to end with "Read
% along any row. The same intervention ... moves in both directions", which is
% the Introduction's own result 3 restated. A caption says what a table shows,
% the body says what it means, and the body already does. Do not grow either
% one back.
%
% The caption covers BOTH kinds of number. It used to read "paired change ...
% against each cell's own baseline", which stopped describing the table the
% moment the baseline row went in, since the first row is an absolute rate and
% not a change. It also names the unit once, as percentage points, so the
% prose can go on saying "points" without either being loose.
%
% The note says "entries" and "column", never "cell". The Introduction uses
% cell for a backbone-benchmark pair, and inside a table the same word means
% a box, so "marks the cell whose checkpoint is not public" pointed at one
% grid cell while eight boxes carried the dash. "Column" is what the reader
% sees and "entry" is what bold marks, and neither collides with the term the
% body has already defined.
%
% The caption covers BOTH kinds of number. It used to say "paired change ...
% against each cell's own baseline", which stopped describing the table the
% moment the baseline row went in: the first row is an absolute rate, not a
% change. It also names the unit once, as percentage points, so the prose can
% keep saying "points" without either being loose.
% ---------------------------------------------------------------------------
\begin{table*}[t]
\centering
\caption{Baseline success rates and paired changes, in percentage points.}
\label{tab:grid}
\setlength{\tabcolsep}{5pt}
\begin{tabular}{l rrr rrr}
\toprule
& \multicolumn{3}{c}{WidowX-Bridge} & \multicolumn{3}{c}{Google Robot/Fractal} \\
\cmidrule(lr){2-4} \cmidrule(lr){5-7}
\multicolumn{1}{c}{Condition} & OpenVLA & SpatialVLA & UniVLA & OpenVLA & SpatialVLA & UniVLA \\
\midrule
Baseline success (\%) & $15.6$ & $30.2$ & $81.2$ & $38.5$ & $84.4$ & \textemdash \\
\midrule
Action repetition 2 & $-8.3$ & $+12.5$ & $\mathbf{-69.8}$ & $+5.2$ & $\phantom{+}0.0$ & \textemdash \\
Action repetition 4 & $\mathbf{-11.5}$ & $-12.5$ & $\mathbf{-81.2}$ & $-1.5$ & $\mathbf{-40.0}$ & \textemdash \\
\addlinespace[2pt]
Visual foveation, log-polar & $+18.8$ & $-8.3$ & $+5.2$ & $\mathbf{-19.3}$ & $+0.7$ & \textemdash \\
Visual foveation, blur & $+17.7$ & $\phantom{+}0.0$ & $-8.3$ & $-8.9$ & $-1.5$ & \textemdash \\
\addlinespace[2pt]
Depth pruning 1 & $+2.1$ & $-10.4$ & $-3.1$ & $+0.7$ & $+8.1$ & \textemdash \\
Depth pruning 2 & $\phantom{+}0.0$ & $-9.4$ & $-4.2$ & $\phantom{+}0.0$ & $+3.0$ & \textemdash \\
Depth pruning 4 & $+1.0$ & $\mathbf{-28.1}$ & $-2.1$ & $\mathbf{+15.6}$ & $\mathbf{-17.8}$ & \textemdash \\
\bottomrule
\end{tabular}

\vspace{3pt}
\parbox{\textwidth}{\footnotesize Every row below the baseline is the paired
change against it, over the episodes present in both runs. \textbf{Bold} marks
the 8 entries that clear a Bonferroni threshold of $\alpha = 0.05/38$ over the
paired tests the grid runs. \textemdash{} marks the column whose checkpoint is
not public.}
\end{table*}
